\pdfmapfile{+cm.map}
\pdfmapfile{+cmextra.map}
\documentclass[10pt,letterpaper]{article}
\usepackage[margin=.60in,headheight=14pt,headsep=13pt,footskip=23pt]{geometry}
\usepackage{fancyhdr,array,amsmath,amssymb,url}
\renewcommand{\familydefault}{\sfdefault}
\pagestyle{fancy}\fancyhf{}
\fancyhead[L]{\fontsize{10}{12}\selectfont Week 77 / Persistent holes / Adult guide}
\fancyfoot[L]{\fontsize{9}{10}\selectfont Bellingham Math Circle / Week 77 / F77-FAC-v1}
\fancyfoot[R]{\fontsize{9}{10}\selectfont\thepage}
\renewcommand{\headrulewidth}{0pt}\renewcommand{\footrulewidth}{0pt}
\setlength{\parindent}{0pt}\setlength{\parskip}{4pt}
\setlength{\tabcolsep}{5pt}\renewcommand{\arraystretch}{1.16}
\raggedright
\newcommand{\sectiontitle}[1]{\par\vspace{5pt}{\bfseries\fontsize{12}{14}\selectfont #1}\par}
\newcommand{\itemtitle}[1]{\par\vspace{3pt}\textbf{#1}\ }
\newcommand{\hint}[1]{\par\textbf{Hints, in order.} #1}
\newcommand{\rim}[1]{\ensuremath{\mathrm{#1}}}
\begin{document}
\sectiontitle{Mathematical destination: facts and limits}
Work in a fixed finite planar triangulation. Triangle interiors do not overlap; pieces meet only along common edges or vertices. Add edges and filled faces without removing anything. Every face needs all three edges; a tied stage is recorded after its whole batch. All cancellation is modulo 2: two copies of one edge cancel.

\textbf{Boundary theorem.} A saved loop is gone precisely when some available filled triangles have its edge set as their combined boundary. Once gone, the same saved loop stays gone in every later growing board. In a planar board a filling, if it exists, is unique; the loop disappears when its last required triangle is filled. A loop's first appearance and final filling do \emph{not} by themselves identify a barcode interval for an arbitrary loop.

\textbf{Counts versus history.} The independent hole count is \(E-V+C-F\), where \(F\) counts filled triangles and \(C\) counts connected components. Counts alone do not determine survival. On the divided square write
\(P=AB+BC+AC,\ Q=AC+CD+DA,\ R=AB+BC+CD+DA\).
Then \(R=P+Q\). Either one-tile filling leaves \(R\) nonzero; both kill it. Two triangles meeting at a single vertex can instead lose their older loop first. Problems 2 and 3 make this distinction explicit.

\textbf{Different drawings, same class.} In Problem 6 four different edge loops represent the same nonzero class after \(ABO,BCO\) are filled. On a fresh unfilled board, their required tile sets differ. Exactly one pair permits either strict disappearance order. The proof uses which tiles are necessary, not guesses about which visible cavity inherits a name.

Children are meant to build, retain an unchanged edge record, test cancellation and find a construction or obstruction. A good first visit reaches one convincing square test and two contrasting histories. Adults support the completeness arguments in Problems 2 and 4 and the general reasoning in Problems 5 and 6. Formal homology, interval decomposition and general stability are adult extensions. A finite check is not an infinite proof; long-lived need not mean meaningful, or short-lived mean noise.

\sectiontitle{Who uses this prototype}
One shared Grades 4--5 packet, five student pages and six problems. Prerequisites: follow short edge labels, distinguish an edge from a filled triangle, keep a saved list unchanged, and cancel matching pairs with support. No multiplication or formal algebra is needed for the first route.

For the stated group, use one kit with the three older children and the organizer: two build/test while the third referees, then rotate each build. The other two adults remain at their fixed K--1 and 2--3 tables with separately prepared activities. This packet supplies no route or print allocation for those eight children. Arrange their activities before using this prototype.

\sectiontitle{Print and prepare: exact target-table allocation}
\textbf{Per pair or this trio:} one shared student packet (pages 1--5, five sheets), handed out one page at a time; one extra copy each of pages 1, 3 and 5 on opaque paper/card for cutting. These yield eight tiles: square \(ABC,ACD\); joined board \(ABC,CDE\); fan \(ABO,BCO,CDO,DAO\). Keep board-specific tiles in three envelopes; identically named triangles have different sizes.

\textbf{For the target table including spares:} print another cutting set of pages 1, 3 and 5 (eight spare tiles) and one spare student packet. Total student printing: \(5+3+3+5=16\) sheets, single-sided US Letter at 100\%. Print one copy of this four-page guide for the organizer. Other tables: zero copies of this packet.

Cut seven yarn/edge strips: four matching the square's 6.2\,cm sides, one its diagonal (about 8.8\,cm), and two matching \(AB,DE\) on page 3 (5.8\,cm). Add two 10\,cm spare lengths to trim. The fan starts with all eight printed edges; no new strips are needed. Prepare the remaining small items listed on the next page.
\newpage
\sectiontitle{Preparation continued; the hour}
\textbf{Per kit:} 33 edge tokens, three each of \(AB,BC,CD,DA,AC,CE,DE,AO,BO,CO,DO\); three blank spare tokens; six stage cards \(0,2,3,4,5,8\); two blank spare stage cards; two contrasting colored pencils. Write labels in either direction consistently. Three copies allow an edge from the saved loop and both incident tile boundaries to be on the table at once.

\textbf{For this table:} add five launch tokens (two \(XY\), two \(YZ\), one \(XZ\)), three more blank token spares, four pencils (one per child plus one spare), three envelopes, one pair of scissors and a little removable tape. Thus there are 44 token cards including launch and blanks. Use the small \(XYZ\) picture on the spare page 1 for the launch. Cut a matching \(XYZ\) tile from a cutting copy's small example; this makes 17 tiles in all, including spares and launch.

\textbf{Before children arrive:} cut both tile sets from matching board copies. Trace/cut rather than resize; mark each tile's name. Keep vertex labels and edge strips visible when placing tiles. On page 5 physically place \(ABO\) and \(BCO\) before starting: the printed board itself has no shaded fillings. Check fit, label visibility and that tokens can be moved without disturbing the saved list or board. Rehearse one whole cancellation and one reset with the real materials. \textbf{Physical rehearsal and classroom piloting have not been performed.}

\textbf{0--5 minutes:} running around. \textbf{5--8:} short whole-group demonstration; only the target trio continues this prototype. Say, \lq\lq An edge stays once it is added. A tile needs all three edges.'' Point to \(XY,YZ,XZ\); ask one child to check them and place the opaque \(XYZ\) tile. Set out \(XY,YZ\) test tokens, add \(XY,YZ,XZ\), and let the child remove equal pairs. Say, \lq\lq Only \(XZ\) remains. We changed test tokens; the board and the saved list stay.'' This is the non-loop example printed on page 1.

\textbf{8--10:} let the trio handle pieces and repeat the move. \textbf{10--20:} Problem 1.
\textbf{20--33:} Problem 2. \textbf{33--45:} Problem 3.
\textbf{45--53:} choose Problem 4 or a short Problem 5 conversation; do not rush both.
\textbf{53--57:} share one saved loop and the evidence for its survival or death.
\textbf{57--60:} pack and record. Problem 6 merits a separate 15--25-minute continuation for a ready group.

\textbf{Stopping/fallback:} if labels or cancellation are still difficult, stay with Problem 1; skip full enumeration, stage changes and the fan. One child chooses a square loop, another predicts whether a single legal tile kills it, and the referee checks with tokens. Swap roles. Ask at sharing: \lq\lq What stayed the same in two builds, and what changed about the saved loop?''

\sectiontitle{Problem 1, page 1: one filling or two? \quad Core}
Children build from dots, save a loop, and test fillings on separate builds.
A one-tile example is \(P=AB,BC,AC\), killed by \(ABC\), or \(Q=AC,CD,DA\), killed by \(ACD\).
The loop surviving either single filling is \(R=AB,BC,CD,DA\).
Add \(AC\) before filling a half, even though it is not on the saved rim.

\textbf{Check with objects:} \(R+P=Q\) and \(R+Q=P\), so one tile cannot remove every token. Using both leaves none. These are the square's only three nonempty loops, up to starting point or direction. No new name for a visible cavity is needed.
\hint{1. Trace one tile's edge places. 2. Keep your saved list and try the other tile on a reset board. 3. Which edges remain after the doubled diagonal is removed?}

\sectiontitle{Problem 2, page 2: every schedule \quad Core}
Children choose both the edge order and the tile order. The complete four-case list is on the next page; exactly three cases meet the target. Keep the first loop's original edge list throughout each replay. If enumeration stalls, find one successful build and move to Problem 3; leave completeness for a return visit.
\newpage
\sectiontitle{Problem 2 continued: complete answer and explanation}
All rows start with \(AB,BC,CD\) at stage 0. A face entry means \lq\lq fill that triangle.''
\begin{center}
\begin{tabular}{ccccccl}
Stage 2 & Stage 4 & Stage 5 & Stage 8 & Saved loop & Gone at & Meets target\\
\hline
\(DA\)&\(AC\)&\(ABC\)&\(ACD\)&\(R\)&8&yes\\
\(DA\)&\(AC\)&\(ACD\)&\(ABC\)&\(R\)&8&yes\\
\(AC\)&\(DA\)&\(ABC\)&\(ACD\)&\(P\)&5&no\\
\(AC\)&\(DA\)&\(ACD\)&\(ABC\)&\(P\)&8&yes
\end{tabular}
\end{center}
Two choices for the first edge and two for the first tile give four cases, with everything else forced. If \(DA\) arrives first, both tiles are needed to fill its saved outer rim. If \(AC\) arrives first, the saved loop is exactly \(ABC\)'s boundary, so only that tile is needed. All four rows have hole counts \(0,1,2,1,0\) at stages \(0,2,4,5,8\).
\hint{1. Which edge did you add first? Save its completed loop before adding anything else. 2. Hold the edge order fixed and switch the two tiles. 3. After both tile choices, switch the edge order. Have you covered both choices at each decision?}

\sectiontitle{Problem 3, page 3: matching counts, different histories \quad Core}
Start with \(AC,BC,CD,CE\). The triangles are \(ABC\) and \(CDE\), meeting only at \(C\). Here is every schedule; the first two rows supply the requested contrast.
\begin{center}
\begin{tabular}{ccccc}
Stage 2 & Stage 4 & Stage 5 & Stage 8 & First loop gone at\\
\hline
\(AB\)&\(DE\)&\(ABC\)&\(CDE\)&5\\
\(AB\)&\(DE\)&\(CDE\)&\(ABC\)&8\\
\(DE\)&\(AB\)&\(ABC\)&\(CDE\)&8\\
\(DE\)&\(AB\)&\(CDE\)&\(ABC\)&5
\end{tabular}
\end{center}
If \(AB\) is first, save \(AB,BC,AC\); if \(DE\) is first, save \(CD,DE,CE\).
Each saved rim disappears exactly when its own tile arrives. The other tile has different edges, so cannot cancel it. \textbf{No}, the identical counts do not tell when the first loop disappears.
\hint{1. Touch the triangle whose missing edge you closed first. 2. Keep that edge list while filling the other triangle first. 3. Can you exchange the tile times without changing any hole count?}

\sectiontitle{Problem 4, page 4: simultaneous arrivals \quad Further}
The printed schedule never has two holes at a finished stage. Its counts at \(0,2,4,8\) are \(0,1,1,0\): the edge and tile at 4 are a single recorded batch. There are four proposed one-item changes:
\begin{center}
\begin{tabular}{p{1.15in}p{.55in}p{4.7in}}
Move & To & Result\\
\hline
edge \(AC\)&3&legal; two holes after stage 3, before the tile at 4\\
edge \(AC\)&5&illegal; \(ABC\) at 4 would lack \(AC\)\\
tile \(ABC\)&3&illegal; its edge \(AC\) would still be absent\\
tile \(ABC\)&5&legal; two holes after stage 4, before the tile at 5
\end{tabular}
\end{center}
Thus exactly the first and fourth changes work. The rim still disappears at 8.
\hint{1. Freeze the board only after all of stage 4 is present. 2. Try moving the edge alone and check every tile's three sides. 3. Now move the tile alone. There are only two new times for each item.}

\sectiontitle{Problem 5, page 4: can a saved loop return? \quad Further}
\textbf{No.} Keep the particular selection of filled triangles that erased every token. All those triangles remain available at every later stage. Reusing exactly that selection still erases exactly the same saved edge list, even if extra edges and tiles have appeared. This proves the claim for every legal growing board; it does not depend on testing examples or on planarity.
\hint{1. Save the tiles you used when the loop disappeared. 2. Which of those tiles could be missing later? 3. Can you replay the identical token cancellation?}
A fresh reset is a new build; it does not contradict this statement.
\newpage
\sectiontitle{Problem 6, page 5: two fixed loops \quad Return visit}
Put \(ABO,BCO\) tiles on the all-edge fan. A loop survives either next filling precisely when its required tile set includes \emph{both} \(CDO\) and \(DAO\). All four possible edge lists follow. Use any two distinct rows for the first request.
\begin{center}
\begin{tabular}{cll}
Label & Saved edges & Required tiles\\
\hline
\(L_0\)&\(AO,CO,CD,DA\)&\(CDO,DAO\)\\
\(L_1\)&\(AB,BO,CO,CD,DA\)&\(ABO,CDO,DAO\)\\
\(L_2\)&\(AO,BO,BC,CD,DA\)&\(BCO,CDO,DAO\)\\
\(L_3\)&\(AB,BC,CD,DA\)&\(ABO,BCO,CDO,DAO\)
\end{tabular}
\end{center}
In this starting board each row needs both remaining tiles, so it survives either first addition and disappears at the last one. The drawings are different; they represent the same current class because their differences are boundaries of already-filled \(ABO,BCO\).

\textbf{For the fresh-board request choose \(L_1,L_2\).}
Fill \(CDO,DAO,ABO,BCO\): \(L_1\) disappears at the third filling and \(L_2\) at the fourth.
Fill \(CDO,DAO,BCO,ABO\): \(L_2\) disappears at the third filling and \(L_1\) at the fourth. Their saved edge lists never change.

\textbf{Why this is the only reversible pair.} Each fan triangle has its own outside edge, appearing in no other tile. Hence the outside edges of a boundary force exactly which tiles were used; no alternative filling can erase a missing required tile. Including both unfilled tiles leaves only the four choices above for the two already-filled tiles. If one required set contains another, the smaller set is completed no later than the larger, in every order. Of the six pairs, only \(L_1,L_2\) have neither set containing the other.
\hint{1. Choose a loop and try each possible first filling on reset boards. 2. Which tiles are inside that loop? Can a missing one be replaced by tiles outside it? 3. On the fresh board, look for two loops that each require a tile the other does not.}
If children can give the two orders but cannot yet show completeness, keep that as a checked construction. The necessity argument is an adult-supported next step.

\sectiontitle{Adult mathematics: short proofs and next steps}
For a finite planar straight-line triangulation, the boundary map on face selections is injective. A nonempty finite union of selected triangles has an exposed edge on its outside boundary, which occurs once. Thus it cannot have zero mod-2 boundary. Two different fillings of the same loop would differ by a nonempty zero-boundary selection, a contradiction. The cycle space of a graph has dimension \(E-V+C\) (use a spanning forest); the \(F\) triangle boundaries are independent. Quotienting by them gives \(E-V+C-F\).

For the square with outer rim born at \(b\), diagonal at \(s\), and faces at \(u,v\), assume \(b<s\leq\min(u,v)\). The old class \(R=P+Q\) survives the first face and vanishes at the second. Choosing the first-filled triangle's rim as the additional basis vector gives intervals
\[
[b,\max(u,v))\quad\hbox{and}\quad[s,\min(u,v)).
\]
Omit a zero-length interval; at a death time the class is already gone. For Problem 4 the second interval is omitted in the original schedule, and becomes \([3,4)\) or \([4,5)\) in the two valid changes. On the joined board, filling older/newer first gives respectively \([2,5),[4,8)\) or \([2,8),[4,5)\). These conclusions use the actual unchanged loop classes across stages, not only hole counts. General persistence pairing is a later adult topic; not every saved loop is an interval generator.

\sectiontitle{Source lineage and what to record}
Edelsbrunner, Letscher and Zomorodian, \emph{Topological Persistence and Simplification} (2002), Section 2 (PDF pp.\ 2--3), Section 3 (pp.\ 4--6): mod-2 boundaries, growing complexes and persistence.
{\footnotesize\url{https://pub.ista.ac.at/~edels/Papers/2002-J-04-TopologicalPersistence.pdf}}
The square, joined-board and fan tasks and the elementary proofs here are independently devised teaching examples. No borrowed figures or source-paper exercises are included.

Afterward record which problems were reached, which moves children controlled, what required adult rescue, and one exact edge list, drawing or claim worth keeping. Mark each claim \emph{tried}, \emph{guessed}, \emph{checked in cases}, or \emph{explained}. Note material-fit and timing problems. Digital verification establishes the printed finite answers; no observed learning or classroom readiness is claimed.
\end{document}
